%══════════════════════════════════════════════════════════════════════════════
%─── Section 4: Theory ───────────────────────────────────────────────────────
%══════════════════════════════════════════════════════════════════════════════
\section{Coverage under Temporal Dependence}
\label{sec:theory}

The split-conformal shift of Section~\ref{sec:cqr_var} carries an exact
finite-sample guarantee when the nonconformity scores are exchangeable
\citep{vovk2005algorithmic,lei2018distribution}. Financial return series are not
exchangeable: volatility clusters, tails arrive together, and the score process
inherits both. This section replaces the exchangeability assumption with a weak
dependence condition and states what survives.

The coverage statement this requires is not new, and the section is organised
around that fact. \citet{oliveira2024split} prove marginal coverage for split
conformal prediction on stationary $\beta$-mixing data with an explicit
remainder, and name GARCH processes as inside the class their result covers;
\citet{barber2026predictive} sharpen the sample-size term from
$\mathcal O(\sqrt{\tau/n})$ to $\mathcal O(\tau/n)$. What this section
contributes is not the bound but its calibration and its price: the mixing rate
of a GARCH score process is made explicit, which turns an abstract separation
requirement into a number of trading days, and
Section~\ref{sec:theory_gap} then measures what imposing that number costs on
the panel this paper reports. That measurement is the part no prior result
supplies, and Section~\ref{sec:blindspot_identification} carries the theoretical claim
this paper does originate.

Three results are stated. The first is asymptotic coverage for the single-split
estimator under geometric $\beta$-mixing, with an explicit remainder and a
corollary that makes the mixing rate concrete for a GARCH data-generating
process. The second is a deterministic bound for an online quantile update, which
holds pathwise and assumes nothing about the data at all. The third is a
finite-window bound for the rolling estimator under distributional drift. They
are stated in that order because the strength of the conclusion falls in that
order, and the last two describe estimators adjacent to, rather than identical
to, the rolling variant this paper runs.

One condition separates the theory from the protocol, and
Section~\ref{sec:theory_gap} states it rather than deferring it.

\subsection{Assumptions}
\label{sec:theory_assumptions}

\begin{assumption}
	\label{ass:main}
	Let $S_t = \hat q_t^{lo} - r_t$ be the nonconformity score process of
	equation~\eqref{eq:score} and let $n = \lfloor f_c T\rfloor$ be the
	calibration sample size.
	\begin{enumerate}[label=\textup{(A\arabic*)},leftmargin=2.6em]
		\item\label{ass:stationary} $\{S_t\}$ is strictly stationary and
		geometrically $\beta$-mixing: $\beta(k)\le C\rho^{k}$ for some
		$\rho\in(0,1)$.
		\item\label{ass:density} The marginal law $F_S$ admits a density $f_S$
		that is continuous and bounded away from zero in a neighbourhood of the
		$(1-\alpha)$-quantile $q_{1-\alpha}$.
		\item\label{ass:weak_dep} There is a gap sequence $g_n\to\infty$ with
		$g_n/n\to 0$ such that the evaluated score $S_{n+g_n+1}$ is separated
		from the calibration block $\mathcal D_n = (S_1,\dots,S_n)$ by $g_n$
		unused observations.
	\end{enumerate}
\end{assumption}

\ref{ass:stationary} is a condition on the score process, not on returns, and it
is not assumed: Lemma~\ref{lem:score_mixing} derives it from the corresponding
condition on returns together with a restriction on the forecaster.
\ref{ass:density} rules out a flat or atomic score distribution at the level
being read; it is the condition under which a quantile is identified at a rate.
\ref{ass:weak_dep} is the separation condition, and Section~\ref{sec:theory_gap}
returns to it.

\subsection{The score process inherits weak dependence, if the context is finite}
\label{sec:theory_lemmas}

\begin{lemma}[Inheritance of weak dependence by the score process]
	\label{lem:score_mixing}
	Let $\{r_t\}$ be strictly stationary and $\beta$-mixing with coefficients
	$\beta_r(k)$. Suppose $\hat q_t^{lo}$ is $\mathcal F_{t-1}$-measurable and
	generated by a time-invariant measurable functional $g$ of a finite context
	of length $m$,
	\begin{equation}
		\label{eq:forecast_functional}
		\hat q_t^{lo} = g(r_{t-1},\dots,r_{t-m}).
	\end{equation}
	Then $S_t = \hat q_t^{lo} - r_t$ is strictly stationary and $\beta$-mixing
	with
	\begin{equation}
		\label{eq:score_mixing_bound}
		\beta_S(k)\;\le\;\beta_r(k-m),\qquad k>m.
	\end{equation}
\end{lemma}

The proof is in Supplement~S.9.1. It turns on a single inclusion,
$\sigma(S_s: s\ge t+k)\subseteq\sigma(r_u: u\ge t+k-m)$, which is what allows a
blocking argument to separate the two $\sigma$-fields.

\begin{remark}[Why the context must be finite, and what EWMA does to it]
	\label{rem:finite_context}
	The finiteness of $m$ is load-bearing. With an infinite context,
	$\hat q_t^{lo} = g(r_{t-1}, r_{t-2},\dots)$, the score $S_s$ depends on
	returns arbitrarily far before $t+k-1$, the inclusion above fails, and no
	separation is obtained --- the bound $\beta_S(k)\le\beta_r(k-1)$ does not
	follow. Restricting $g$ to a window of length $m$ is what makes the argument
	available.

	For all but one forecaster in this study the restriction holds by
	construction: $m = 512$ for the foundation models and $m = 250$ for the
	GARCH-family and historical-simulation benchmarks
	(Section~\ref{sec:exp_design}). EWMA is the exception. Its conditional
	variance is the RiskMetrics recursion, whose memory is infinite, so
	Lemma~\ref{lem:score_mixing} does not apply to it as literally implemented.
	The quantitative position is that at $\lambda = 0.94$ the weight beyond lag
	$k$ is $\lambda^{k}$, which is $\nLitLambdaTail$ at $k = 250$ and falls below
	double-precision machine epsilon at $k = 583$; the lemma applies to the filter
	truncated at $m = 600$, which is indistinguishable from the recursion in the
	arithmetic that produced the series. That is an approximation argument rather
	than an exemption, and it is the only one in this paper.
\end{remark}

\begin{lemma}[Empirical quantile deviation under $\beta$-mixing]
	\label{lem:quantile_mixing}
	Under Assumption~\ref{ass:main}, with $\qVstat$ the calibration-sample
	conformal shift of equation~\eqref{eq:qV} and
	$q_{1-\alpha} = F_S^{-1}(1-\alpha)$,
	\begin{equation}
		\label{eq:quantile_rate}
		\left|\qVstat - q_{1-\alpha}\right|
		=
		\mathcal O_p\!\left(
		\frac{1}{f_S(q_{1-\alpha})}\sqrt{\frac{\log n}{n}}
		\right).
	\end{equation}
\end{lemma}

The rate comes from the uniform concentration of the empirical distribution
function under $\beta$-mixing \citep[Theorem~2.1]{yu1994rates}
\citep[Theorem~7.2]{rio2017asymptotic}, transferred to the quantile scale by
\ref{ass:density}. The $\sqrt{\log n}$ is the price of dependence: under
independence the same statement holds without it.

Two conventions differ at the order the lemma controls, and the distinction is
recorded here because it recurs. Equation~\eqref{eq:qV} takes the
$\lceil (n+1)(1-\alpha)\rceil$-th order statistic; the plain empirical quantile
$Q_{1-\alpha}$ takes one order statistic lower. The gap is $\mathcal O(1/n)$ and
is absorbed into the remainder of Theorem~\ref{thm:dependent_conformal_var}
without changing its rate. It is not absorbed in finite samples, where it has a
sign: at $\alpha = 0.01$ and $n = \nTheoryNCal$ the conformal index sits at
$k/n = \nTheoryConfIndex$ rather than $0.99$, so $\qVstat$ exceeds the sample
$(1-\alpha)$-quantile, and Section~\ref{sec:mc} measures what that costs a
correctly specified forecaster.

\subsection{Asymptotic coverage for the single-split estimator}
\label{sec:theory_theorem}

\begin{theorem}[Coverage under geometric $\beta$-mixing]
	\label{thm:dependent_conformal_var}
	Under Assumption~\ref{ass:main}, the single-split conformal Value-at-Risk
	estimator with $\qVstat$ given by equation~\eqref{eq:qV} satisfies
	\begin{equation}
		\label{eq:coverage_bound}
		\P\!\left(
		r_{n+g_n+1}\;\ge\;\hat q^{lo}_{n+g_n+1}-\qVstat
		\right)
		\;\ge\;
		1-\alpha-\Delta_n,
		\qquad
		\Delta_n \;\le\; C\sqrt{\frac{\log n}{n}}+\beta(g_n),
	\end{equation}
	where $C>0$ depends on the score density near $q_{1-\alpha}$ and on the
	mixing rate.
\end{theorem}

\paragraph{Attribution.} Theorem~\ref{thm:dependent_conformal_var} is a
one-sided specialization of a published result and is stated here because the
paper's estimator needs it, not because it is new. \citet[Theorem~4]{oliveira2024split}
give marginal coverage for split conformal prediction under stationary
$\beta$-mixing with remainder $\varepsilon_{\mathrm{cal}} +
\varepsilon_{\mathrm{train}} + \delta_{\mathrm{cal}}$, in which
$\varepsilon_{\mathrm{cal}}$ carries a $\sqrt{\log(\cdot)/n_{\mathrm{cal}}}$
term and $\varepsilon_{\mathrm{train}}$ is a $\beta$-coefficient evaluated at a
separation; their statement is for a generic score, so a signed score is admitted
without modification. \citet{barber2026predictive} then improve the sample-size
term to scale linearly in $\tau/n$ rather than as $\sqrt{\tau/n}$, which under
geometric mixing with $\tau = c\log n$ is $\mathcal O((\log n)/n)$ against the
$\mathcal O(\sqrt{(\log n)/n})$ of equation~\eqref{eq:coverage_bound}. A reader
who wants the sharpest available constant should take theirs. Equation~\eqref{eq:coverage_bound}
is retained in the form above because it is the form the panel is measured
against in Section~\ref{sec:theory_gap}, and because its two terms separate the
two things a practitioner controls.

The proof, in Supplement~S.9.3, has three steps: the coverage event is rewritten
as a score event, $\qVstat$ is placed near $q_{1-\alpha}$ by
Lemma~\ref{lem:quantile_mixing}, and the dependence between the calibration block
and the evaluated point is bounded by $\beta(g_n)$. The two terms of $\Delta_n$
are separately interpretable and separately controllable: $C\sqrt{(\log n)/n}$ is
estimation error in the calibration quantile and falls with the calibration
sample; $\beta(g_n)$ is residual dependence and falls with the separation.

The conclusion is asymptotic. No finite-sample coverage claim is made anywhere in
this paper for the dependent case, and the bound is conservative at financial
sample sizes: on \nBoundPairs{} representative pairs spanning score persistence
$\hat\rho\in[\nBoundRhoLo,\nBoundRhoHi]$ the estimated remainder lies between
$\nBoundDeltaLo$ and $\nBoundDeltaHi$, a guaranteed coverage floor of
\nBoundFloor\% against empirical coverage averaging \nBoundEmpirical\%. The bound justifies the
correction under dependence; it does not certify it.

\begin{corollary}[GARCH data-generating processes]
	\label{cor:garch}
	Let $r_t = \sigma_t\varepsilon_t$ follow a GARCH$(1,1)$ recursion
	$\sigma_t^2 = \omega + \alpha_1 r_{t-1}^2 + \beta_1\sigma_{t-1}^2$ with
	$\omega>0$, $\alpha_1,\beta_1\ge 0$ and
	$\E[\log(\alpha_1\varepsilon_t^2+\beta_1)]<0$, and let $\varepsilon_t$ have a
	continuous density that is positive near the $\alpha$-quantile. Let the
	forecaster satisfy equation~\eqref{eq:forecast_functional} with finite
	context $m$. Then the score process is geometrically $\beta$-mixing, and
	choosing
	\begin{equation}
		\label{eq:garch_gap_rate}
		g_n = c\log n,
		\qquad
		c > 1/\lvert\log\rho\rvert,
	\end{equation}
	gives $\beta(g_n) = \mathcal O(n^{-c\lvert\log\rho\rvert}) = o\!\left(\sqrt{(\log n)/n}\right)$,
	so that Theorem~\ref{thm:dependent_conformal_var} applies with
	$\Delta_n = \mathcal O\!\left(\sqrt{(\log n)/n}\right)$.
\end{corollary}

The corollary is the reason the theorem is usable rather than merely true. Its
hypotheses are the standard conditions for strict stationarity and geometric
ergodicity of GARCH$(1,1)$ \citep[Theorem~2.10]{francq2019garch,lindner2009stationarity},
and they are satisfied by every parametric benchmark in this study. The
logarithmic rate and its constant are not new either: for geometrically ergodic
Markov chains with rate $\rho$, \citet[Theorem~5.1]{zheng2024conformal} obtain
a separation of order $\log n/\lvert\log\rho\rvert$, which is
equation~\eqref{eq:garch_gap_rate}. Their construction is a thinning stride --- a
$K$-split scheme retaining every $K$th calibration score --- rather than a single
block separated from the evaluated point, so the two are the same rate applied to
different objects. What the corollary adds is the passage from a GARCH
specification to that rate, which rests on the exponential $\beta$-mixing of
GARCH$(1,1)$ established by \citet{carrasco2002mixing} and is a citation rather
than a result. Its value is operational: it converts a mixing condition into a
number of trading days, and the number turns out to be small. Under it the
required separation is logarithmic in the calibration sample: across the
\nGapPanelCells{} cells of the panel $g_n$ runs from \nGapPanelGapLo{} to
\nGapPanelGapHi{} observations with a median of \nGapPanelGapMed, so the
condition is cheap where it can be met at all. The proof
is in Supplement~S.9.4.

\subsection{The separation condition, and what the protocol does instead}
\label{sec:theory_gap}

\ref{ass:weak_dep} requires a gap $g_n\to\infty$ between the calibration block
and the evaluated point, and by Lemma~\ref{lem:score_mixing} that gap must exceed
the forecaster's context length $m$ before the mixing bound binds at all. The
protocol as originally run was a contiguous $70/30$ split with $g_n = 0$, under
which Theorem~\ref{thm:dependent_conformal_var} does not apply to the estimator
at all.

This is the standard operational choice in conformal prediction, and the
distance between it and the theorem is measured rather than argued about. The
gapped estimator sets $g_n = \lceil c\log n\rceil$ with
$c = 1/\lvert\log\hat\rho\rvert$, the minimal gap
satisfying~\eqref{eq:garch_gap_rate}, taken out of the front of the test block so
that the shift is computed on the same calibration sample and only the evaluation
window moves. It is computed on every cell of the panel and compared against the
estimator the tables report.

The constant is undefined where $\hat\rho \le 0$, and a floor of five
observations is imposed there. That occurs on \nRhoFallback{} of the
\nRhoCells{} cells, and \nRhoFallbackTrunc{} of those are the two truncated
Chronos series: the median $\hat\rho$ is $\nRhoMedTrunc$ across their
\nRhoTruncCells{} cells against $\nRhoMedOther$ across the remaining
\nRhoOtherCells. The concentration follows from what the truncation does.
Restricting the predictive law to its most probable atoms leaves a lower
quantile that moves little between adjacent days while the return does not, so
the score $\hat q^{lo}_t - r_t$ is dominated by $-r_t$ and inherits the return's
lack of persistence. A degenerate constant is therefore the corollary reporting
that those scores carry no dependence to separate, and the floor supplies more
separation than the condition requires; the fallback is not concentrated where
the gap would bind hardest but where it would bind least. Both readings of
the census --- concentration on the truncated series, and no forecaster pattern
at all --- were written down before it ran, together with the rule that decides
between them, and the second would have been reported. That ordering is not
evidenced the way the two dated declarations elsewhere in this paper are: the
declaration and the output are minutes apart in one commit, so what a reader can
verify is the alternative on the page and not the sequence in which it was
written.

One property of the construction makes the comparison sharper than an ablation
can be. Because the gap is removed from the front of the test block, the
calibration sample is untouched, and the shift is therefore not re-estimated:
across all \nGapAllCells{} cells of the panel the two arms return
\emph{identical} values of $\qVstat$, to the last bit. The two arms are the same
threshold evaluated on different days. Every difference reported below is
residual dependence between the calibration block and the evaluated point and
cannot be estimation noise, because there is no second estimate.

Across all \nGapAllCells{} cells, at \nGapAllLevels{} levels, the gap runs from
\nGapAllGapLo{} to \nGapAllGapHi{} observations, a median of
\nGapAllGapMed{} and at widest \nGapWidestPct\% of the test window. It moves
the corrected violation rate by at most $\nGapAllDPiMax$; it changes
\nGapAllZoneChanges{} Basel zone classifications and flips
\nGapAllKupFlips{} Kupiec verdicts at the $5\%$ level. At $\alpha = 0.01$, the
level every headline table reports, the corresponding counts are
\nGapPanelZoneChanges{} and \nGapPanelKupFlips{} across \nGapPanelCells{}
cells, with the rate moving by a median of $\nGapPanelDPiMed$ and at most
$\nGapPanelDPiMax$. Those flips are what a threshold does to a statistic that
moved by less than a thousandth, not instability in the estimator.

The rate $g_n = \lceil c\log n\rceil$ has a regime in which it cannot be
imposed at all: as $\hat\rho$ approaches one the constant diverges, and the
required separation exceeds the test block. At the sample sizes here that
happens only above $\hat\rho = \nGapRhoRaiseAt$, and the largest value on the
panel is \nGapMaxRho, so the condition is unmet on \nGapCellsUnmet{} cells and
the widest gap consumes \nGapWidestPct\% of its evaluation window. The
implementation raises on such a cell rather than truncating the gap, because a
coverage figure computed on whatever observations survived a silent truncation
is not a coverage figure. Should a series ever reach that regime, it is reported
as a cell on which \ref{ass:weak_dep} cannot be met at this sample size, not
excluded and not covered.

So Theorem~\ref{thm:dependent_conformal_var} covers an estimator that differs
from the one the tables report by at most $\nGapAllDPiMax$ in violation rate,
with \nGapAllZoneChanges{} regulatory classifications and
\nGapAllKupFlips{} Kupiec verdicts changed across \nGapAllCells{} cells. That
is a measured distance and not a proof that the theorem applies to the
contiguous split, which it does not. The results below are reported on the
contiguous split, as they were computed; a reader who requires the covered
estimator can read the difference off the figures above. What the theorem does
not cover at all is the rolling variant, for the separate reason given in
Section~\ref{sec:theory_rolling}.

\subsection{The rolling estimator: two bounds, and what they cover}
\label{sec:theory_rolling}

The rolling variant of equation~\eqref{eq:qV_roll} recomputes the shift from a
trailing window, which violates the single-split assumption
Theorem~\ref{thm:dependent_conformal_var} rests on. Two results bound related
objects. Neither covers the uniform-weight trailing-window quantile exactly, and
the scope of each is stated with it.

\begin{proposition}[Time-average coverage of the online quantile update]
	\label{prop:ogd_regret}
	Let $q_1\in[-R,R]$ evolve by
	\begin{equation}
		\label{eq:ogd_update}
		q_{t+1} = q_t + \eta\left[\mathbf 1\{s_t>q_t\}-\alpha\right],
	\end{equation}
	with step size $\eta>0$ and scores satisfying $|s_t|\le R$. Then
	\begin{equation}
		\label{eq:ogd_bound}
		\left|\frac1T\sum_{t=1}^{T}\mathbf 1\{s_t>q_t\}-\alpha\right|
		\;\le\;
		\frac{2R+\eta}{\eta T},
	\end{equation}
	and with $\eta = R/\sqrt T$ the right-hand side is at most $3/\sqrt T$.
\end{proposition}

The bound is deterministic. It holds for every score sequence, with no
probabilistic assumption whatsoever, and the proof in Supplement~S.9.5 needs only
two observations: the iterates cannot leave a bounded interval, because the
update reverses sign as soon as they try; and the recursion telescopes exactly.
The usual route to a statement of this kind invokes the online-convex-optimisation
regret guarantee of \citet{zinkevich2003online} and converts average regret on
the pinball loss into average miscoverage, a conversion that requires a local
Lipschitz argument in the indicator parametrisation and does not follow from the
regret bound alone.

Three limitations bound its reach. It controls the time average of the coverage
indicator, not conditional coverage and not coverage over any subwindow: a
sequence can satisfy~\eqref{eq:ogd_bound} while violating badly in one regime and
over-covering in another. It requires bounded scores. And it is a statement about
the online update~\eqref{eq:ogd_update}, whereas the estimator run here is a
trailing-window empirical quantile. In stationary regimes the two target the same
population quantile; under drift the window trades bias for variance near
transitions.

\begin{proposition}[Finite-window coverage under drift]
	\label{prop:rolling_drift}
	Let $\qVroll$ be the trailing-window shift of
	equation~\eqref{eq:qV_roll} over a window of size $w$, and let
	\begin{equation}
		\label{eq:drift_measure}
		\delta_w(t) = \max_{\tau\in[t-w,\,t-1]}
		d_{\mathrm{TV}}\!\left(\mathcal L(S_\tau),\mathcal L(S_t)\right).
	\end{equation}
	Under local stationarity conditions analogous to
	Assumption~\ref{ass:main},
	\begin{equation}
		\label{eq:rolling_bound}
		\P\!\left(r_t\ge\hat q_t^{lo}-\qVroll\right)
		\;\ge\;
		1-\alpha-\delta_w(t)-\Delta_w,
		\qquad
		\Delta_w\le C\sqrt{\frac{\log w}{w}}.
	\end{equation}
\end{proposition}

The bound separates distributional drift, in $\delta_w(t)$, from finite-sample
estimation error under dependence, in $\Delta_w$. Its conclusion holds under the
total-variation drift measure~\eqref{eq:drift_measure}, but that measure is not
estimable from a single path without further assumptions, and the constant it
produces is too loose to be informative at $\alpha = 0.01$. It is retained
because it makes precise \emph{which} quantity has to be small for a trailing
window to work, and $\delta_w(t)$ is used operationally as a monitoring
diagnostic rather than as a certified bound. The empirical proxy
$\hat\delta_w(t)$, the total variation distance between the first and second
halves of the window over $\sqrt w$ bins, is defined in Supplement~S.9.6 and
plotted for Lag-Llama on the S\&P~500 in Supplement Figure~S.2; it flags the
global financial crisis, the COVID-19 shock and the March 2023 banking episode.

\begin{remark}[What the rolling estimator is entitled to claim]
	\label{rem:rolling_status}
	The formal conformal guarantee, Theorem~\ref{thm:dependent_conformal_var},
	covers the single-split estimator. The rolling estimator is a locally
	adaptive engineering variant whose validity is assessed empirically, and
	supported by the time-average bound of Proposition~\ref{prop:ogd_regret} and
	the drift bound of Proposition~\ref{prop:rolling_drift} --- neither of which
	is the estimator itself. It is described as a heuristic wherever it appears,
	and the empirical comparison against the static variant is in
	Section~\ref{sec:recal} and Supplement~S.12.1.
\end{remark}
